\chapter{Scope, evidence and physical assumptions}
\section{What this document establishes}
This report derives an explicit six-degree-of-freedom model, a six-controller
architecture, the covariance and gain calculations used by the C estimators,
and a bounded local-navigation baseline. Numerical results are generated by the
production C control equations and a common C++ navigation core. They establish
behaviour under stated simulation assumptions. They do not identify the actual
aircraft or certify a flight tune.

The physical aircraft is an F450/Q450-style X quadcopter with 9047 propellers,
nominal 1400 KV motors, four 40 A ESCs and a 3S 4200 mAh LiPo. The approximately
150 g battery mass is unverified. A motor's KV value and an ESC's current rating
do not determine thrust, propeller efficiency or inertia. In particular, 40 A is
not a thrust measurement. The following complete-vehicle model is an explicit
test assumption, not the previous frame-plus-battery mass subtotal.
\begin{center}
\begin{tabular}{lll}\toprule
Quantity & Test value & Unit\\\midrule
Complete mass $m$ & 1.2 & kg\\
Diagonal motor spacing $L$ & 0.45 & m\\
$J=\operatorname{diag}(J_x,J_y,J_z)$ & $(0.012,0.012,0.022)$ & kg m$^2$\\
Maximum thrust per motor $T_{\max}$ & 6 & N\\
Reaction torque coefficient $\kappa$ & 0.015 & m\\
Thrust response time $\tau_m$ & 0.04 & s\\
Linear drag coefficient $d_v$ & 0.12 & N s/m\\
Gravity $g$ & 9.80665 & m/s$^2$\\
Control/physics interval $\Delta$ & 0.002 & s\\
Barometer update rate & 50 & Hz\\
Gyro sample standard deviation & 0.2 & deg/s\\
Accelerometer sample standard deviation & 0.005 & $g$ per axis\\
Barometric height standard deviation & 0.05 & m\\\bottomrule
\end{tabular}
\end{center}
The collision envelope uses $L/2+0.2286/2=0.3393$ m, rounded to
0.34 m. This includes the 9-inch propeller radius. Using $L/2$ alone would
overstate obstacle clearance by approximately 11.5 cm. The actual geometry and
propeller clearance must be measured; the circle is conservative in plan view.

\section{Evidence boundaries}
\begin{enumerate}
\item The CPU plant integrates rigid-body motion, motor lag and sensor noise.
The shared library calls the same C functions as the flight firmware's math
core. Host compilation is not RP2350 instruction emulation or hardware timing.
\item Two Python ROS 2 adapter nodes exchange test envelopes through Fast DDS
UDP and invoke that C library. The C++ planner is compiled as a shared library.
The ROS-specific C++ wrapper requires a full ROS development environment and
has not been compiled here.
\item YOLO11n inference is tested independently on an official reference image.
The original airborne baseline route uses a known static obstacle map and simulated local
odometry, not YOLO-generated obstacle ranges.
\item Isaac Sim launch attempts failed before a healthy simulation server was
available: physical-device enumeration/GPU initialization failed. There is no
successful Isaac physics run or Isaac recording in these results. The delivered
video is an original visualization of logged CPU simulation data.
\item Physical motors remain disabled. No UART, LoRa, physical IMU, ESC, or real
camera flight pipeline is exercised by the tests in this report.
\end{enumerate}

\chapter{Stochastic modelling foundations}
\section{Covariance and its transformation}
Let $x$ have mean $\mu$ and covariance
$P=\E[(x-\mu)(x-\mu)^{\mathsf T}]$. For $y=Ax+b$,
\begin{align}
y-\E[y]&=A(x-\mu),\\
\operatorname{Cov}(y)
 &=\E[A(x-\mu)(x-\mu)^{\mathsf T}A^{\mathsf T}]
 =APA^{\mathsf T}.
\end{align}
For independent zero-mean process noise $w$ of covariance $Q$,
$\operatorname{Cov}(Ax+w)=APA^{\mathsf T}+Q$: the two cross expectations
vanish by independence and zero means. If the noise is correlated with the
state, this last step is invalid and cross-covariance terms must be retained.

Symmetry follows because $P_{ij}=\E[e_i e_j]=P_{ji}$. For every $v$,
$v^{\mathsf T}Pv=\E[(v^{\mathsf T}e)^2]\geq0$, so a covariance is positive
semidefinite. For a symmetric $2\times2$ matrix, necessary and sufficient
conditions are $P_{00}\geq0$, $P_{11}\geq0$ and
$P_{00}P_{11}-P_{01}^2\geq0$. Covariance units vary by entry; for height and
speed, $P_{00}$ is m$^2$, $P_{01}$ is m$^2$/s and $P_{11}$ is m$^2$/s$^2$.

\section{Height kinematics from direct time integration}
If acceleration $a_k$ is constant during one interval, integrate $\dot v=a_k$
from $0$ to $t$ to obtain $v(t)=v_k+a_kt$. Integrate $\dot h=v(t)$ from zero
to $\Delta$: $h_{k+1}=h_k+v_k\Delta+a_k\Delta^2/2$, while
$v_{k+1}=v_k+a_k\Delta$. Thus the altitude state update uses
\begin{equation}
F=\begin{bmatrix}1&\Delta\\0&1\end{bmatrix},\qquad
G=\begin{bmatrix}\Delta^2/2\\\Delta\end{bmatrix}.
\end{equation}
These are the particular physical kinematic equations; no general matrix
multiplication or matrix-exponential tutorial is included.

\section{Two distinct process-noise conventions}
For one independent acceleration error $\eta_k$ held constant for an entire
sample, $\operatorname{Var}(\eta_k)=\sigma_a^2$ and $w_k=G\eta_k$.
Thus
\begin{equation}\label{eq:qheld}
Q_{\rm held}=\sigma_a^2GG^{\mathsf T}
=\sigma_a^2\begin{bmatrix}\Delta^4/4&\Delta^3/2\\\Delta^3/2&\Delta^2\end{bmatrix}.
\end{equation}
For continuous white acceleration with spectral density $q_c$, however,
\begin{align}
Q_{\rm continuous}
&=\int_0^\Delta e^{A\tau}Bq_cB^{\mathsf T}e^{A^{\mathsf T}\tau}\dd\tau\\
&=q_c\int_0^\Delta\begin{bmatrix}\tau^2&\tau\\\tau&1\end{bmatrix}\dd\tau
=q_c\begin{bmatrix}\Delta^3/3&\Delta^2/2\\\Delta^2/2&\Delta\end{bmatrix}.
\end{align}
These are different stochastic models with different units. The altitude C
implementation uses \eqref{eq:qheld}, with $\sigma_a=1.5$ m/s$^2$ per
2 ms sample. This is not a continuous spectral density. Changing the sample
rate while holding that number fixed changes the effective uncertainty per
second. Real noise parameters require measured bandwidth and timing.

\chapter{Coordinate frames and rigid-body mechanics}
\section{Frames, rotations and signs}
World $W$ is a local right-handed frame with $z$ upward. Body $B$ uses forward,
left, up (FLU). $R=R_z(\psi)R_y(\theta)R_x(\phi)$ maps body vectors to world
vectors. Here $\phi,\theta,\psi$ are roll, pitch, yaw in radians for mechanics.
The elementary matrices are
\begin{align}
R_x&=\begin{bmatrix}1&0&0\\0&c_\phi&-s_\phi\\0&s_\phi&c_\phi\end{bmatrix},\quad
R_y=\begin{bmatrix}c_\theta&0&s_\theta\\0&1&0\\-s_\theta&0&c_\theta\end{bmatrix},\\
R_z&=\begin{bmatrix}c_\psi&-s_\psi&0\\s_\psi&c_\psi&0\\0&0&1\end{bmatrix},
\end{align}
where $c_\alpha=\cos\alpha$ and $s_\alpha=\sin\alpha$. Multiplication yields
\begin{equation}
R=\begin{bmatrix}
c_\psi c_\theta&c_\psi s_\theta s_\phi-s_\psi c_\phi&c_\psi s_\theta c_\phi+s_\psi s_\phi\\
s_\psi c_\theta&s_\psi s_\theta s_\phi+c_\psi c_\phi&s_\psi s_\theta c_\phi-c_\psi s_\phi\\
-s_\theta&c_\theta s_\phi&c_\theta c_\phi
\end{bmatrix}.
\end{equation}
For example, entry $(3,3)$ is $c_\theta c_\phi$, so upward thrust contributes
$T c_\theta c_\phi$ to vertical force. At zero yaw and small angles, positive
pitch accelerates forward and positive roll accelerates right, which is
negative world $y$ when body and world headings align. This explains the signs
in the companion acceleration-to-attitude conversion.

Body angular velocity is $\omega=[p,q,r]^{\mathsf T}$. It is not generally the
Euler-angle derivative:
\begin{equation}
\begin{bmatrix}\dot\phi\\\dot\theta\\\dot\psi\end{bmatrix}
=\begin{bmatrix}1&s_\phi\tan\theta&c_\phi\tan\theta\\
0&c_\phi&-s_\phi\\0&s_\phi/c_\theta&c_\phi/c_\theta\end{bmatrix}
\begin{bmatrix}p\\q\\r\end{bmatrix}.
\end{equation}
Near hover $p\simeq\dot\phi$, $q\simeq\dot\theta$, $r\simeq\dot\psi$.
The independent roll/pitch filters and gain design use that approximation;
the plant itself propagates a unit quaternion. Aggressive attitude changes
require a coupled quaternion estimator and appropriate rate-setpoint mapping.

\section{Newton and Euler equations}
Let $p_W$ be position, $v_W$ velocity, and $e_3=[0,0,1]^{\mathsf T}$. Then
\begin{align}
\dot p_W&=v_W,\\
m\dot v_W&=R e_3\sum_{i=1}^4T_i-mge_3-d_vv_W+f_{\rm wind},\\
J\dot\omega&=\tau-\omega\times(J\omega).
\end{align}
The rotational equation follows from angular-momentum balance. If $h_B=J\omega$,
then the inertial derivative expressed in rotating body axes is
$\dot h_B+\omega\times h_B=\tau_B$. For diagonal inertia,
\begin{align}
J_x\dot p&=\tau_x+(J_y-J_z)qr,\\
J_y\dot q&=\tau_y+(J_z-J_x)rp,\\
J_z\dot r&=\tau_z+(J_x-J_y)pq.
\end{align}
The inertia tensor must be positive definite and satisfy physical principal
moment inequalities; the provisional diagonal values do. A measured model
requires component masses and positions, using
\begin{equation}
J=\sum_i\left(R_iJ_{i,\rm COM}R_i^{\mathsf T}
 +m_i[(r_i^{\mathsf T}r_i)I-r_ir_i^{\mathsf T}]\right).
\end{equation}
This is the parallel-axis theorem. It separates each component's own inertia
from the extra moment caused by its distance from the vehicle centre of mass.

\section{Motor geometry, allocation and saturation}
Set $a=L/(2\sqrt2)=0.159099$ m. Logical motor order is front-left (FL),
front-right (FR), rear-right (RR), rear-left (RL), at
$(a,a,0),(a,-a,0),(-a,-a,0),(-a,a,0)$ respectively. For
$r_i=(x_i,y_i,0)$ and thrust $f_i=(0,0,T_i)$,
\begin{equation}
r_i\times f_i=\begin{bmatrix}y_iT_i\\-x_iT_i\\0\end{bmatrix}.
\end{equation}
Assigning alternating reaction-torque signs $(-,+,-,+)$ gives
\begin{equation}
\begin{bmatrix}T\\\tau_x\\\tau_y\\\tau_z\end{bmatrix}
=\begin{bmatrix}1&1&1&1\\a&-a&-a&a\\-a&-a&a&a\\-\kappa&\kappa&-\kappa&\kappa\end{bmatrix}
\begin{bmatrix}T_{FL}\\T_{FR}\\T_{RR}\\T_{RL}\end{bmatrix}.
\end{equation}
These are assumed spin directions; wiring must eventually be checked. In the
software, $u_i$ is a normalized \emph{desired thrust proxy}, not a verified PWM,
DShot, RPM or electrical-power relationship. Define
\begin{equation}
\begin{bmatrix}u_{FL}\\u_{FR}\\u_{RR}\\u_{RL}\end{bmatrix}
=\begin{bmatrix}1&1&-1&-1\\1&-1&-1&1\\1&-1&1&-1\\1&1&1&1\end{bmatrix}
\begin{bmatrix}c\\u_\phi\\u_\theta\\u_\psi\end{bmatrix}.
\end{equation}
Before saturation and lag, $T_i=T_{\max}u_i$, hence
\begin{equation}
T=4T_{\max}c,\quad\tau_x=4aT_{\max}u_\phi,\quad
\tau_y=4aT_{\max}u_\theta,\quad\tau_z=4\kappa T_{\max}u_\psi.
\end{equation}
For example, positive pitch must increase rear thrust and decrease front
thrust. A host regression test now checks this sign. The original allocator clipped outputs independently; the current implementation
scales infeasible differential spans and shifts collective into a feasible interval,
as derived in the boost/authority chapter. Saturation still limits physical authority.

Motor lag is represented by $\tau_m\dot u_i^{\rm actual}+u_i^{\rm actual}=u_i$.
For constant commanded $u_i$ over $\Delta$, solve the scalar differential
equation using an integrating factor:
\begin{equation}
u^{\rm actual}_{k+1}=e^{-\Delta/\tau_m}u^{\rm actual}_k
 +(1-e^{-\Delta/\tau_m})u_k.
\end{equation}
This exponential update is exact for the assumed first-order actuator. Actual
motors may need $T=k_T\Omega^2$, $\tau=k_Q\Omega^2$ and a measured electrical
model. No battery sag or propeller aerodynamics are currently identified.

\section{Numerical propagation}
Position uses $p_{k+1}=p_k+\Delta v_k+\tfrac12\Delta^2a_k$ and velocity uses
$v_{k+1}=v_k+\Delta a_k$. Angular velocity uses a first-order update of the
Euler equation. Given the updated $\omega$ held over a step,
\begin{equation}
\delta q=\begin{bmatrix}\cos(\alpha/2)\\
\hat\omega\sin(\alpha/2)\end{bmatrix},\quad
\alpha=\|\omega\|\Delta,\quad q_{k+1}=q_k\otimes\delta q,
\end{equation}
with the zero-speed limit handled continuously and the quaternion normalized.
For $q=(s,v)$ and $r=(t,w)$, multiplication is
$q\otimes r=(st-v^{\mathsf T}w,\,sw+tv+v\times w)$.
This preserves rotational geometry for the held angular velocity, but the
entire plant integrator is not an exact solution of coupled nonlinear motion.
Timestep refinement and empirical identification are separate requirements.

\chapter{Constructing the flight model step by step}
\section{Start from a free-body diagram, not a guessed transfer function}
Our first modelling choice is a system boundary: the complete flying aircraft,
including battery, frame, motors, electronics and propellers. Its total mass is
$m=\sum_i m_i$, and its centre of mass is
$r_C=\sum_i m_i r_i/m$. Choosing the body origin at $r_C$ removes the extra
moment caused by expressing angular momentum about an accelerating offset point.
Neither a battery mass nor a frame mass is the complete $m$ in Newton's law.

Choose a local world frame with upward unit vector $e_3$ and a forward-left-up
body frame. Let the four thrust forces be $T_i e_3$ in body coordinates.
There are three categories of external force: rotated rotor thrust, weight,
and aerodynamic/disturbance force. Internal frame forces cancel when summing
over the system boundary. Thus the force balance is constructed as
\begin{align}
\sum f_W&=\underbrace{R e_3(T_1+T_2+T_3+T_4)}_{\text{rotors}}
 -\underbrace{mg e_3}_{\text{weight}}
 -\underbrace{d_vv_W}_{\text{linear drag}}+f_w,\\
\dot p_W&=v_W,\qquad m\dot v_W=\sum f_W.
\end{align}
Dividing by $m$ produces acceleration, not velocity. Integrating acceleration
once gives velocity and a second time gives position. This is why a translational
position transfer function has two low-frequency integrations before damping,
control or other dynamics are introduced.

For each motor at $r_i=(x_i,y_i,0)$, the torque about the centre of mass is
$(y_iT_i,-x_iT_i,0)$. Add rotor drag-reaction torques with alternating signs.
The sum is the allocation equation derived in the mechanics chapter. Newton's
angular counterpart is $J\dot\omega+\omega\times J\omega=\tau$. We keep the
cross term in the nonlinear plant and explicitly neglect its second-order
perturbation when constructing the hover linear model.

\section{Write the nonlinear differential equations before linearizing}
Write $T=\sum T_i$ and use the third column of $R$. The translational equations
are, component by component,
\begin{align}
m\ddot x&=T(c_\psi s_\theta c_\phi+s_\psi s_\phi)-d_v\dot x+f_{wx},\\
m\ddot y&=T(s_\psi s_\theta c_\phi-c_\psi s_\phi)-d_v\dot y+f_{wy},\\
m\ddot h&=Tc_\theta c_\phi-mg-d_v\dot h+f_{wz}.
\end{align}
For diagonal inertia,
\begin{align}
\dot p&=[\tau_x+(J_y-J_z)qr]/J_x,\\
\dot q&=[\tau_y+(J_z-J_x)rp]/J_y,\\
\dot r&=[\tau_z+(J_x-J_y)pq]/J_z.
\end{align}
These six acceleration equations need kinematics as well: $\dot p_W=v_W$
and quaternion attitude propagation. The state is position, velocity, attitude
and angular velocity; four additional motor-lag states make the actuator
dynamics explicit. Quaternion norm is constrained to one, so four quaternion
coordinates represent three independent orientation degrees of freedom.

\section{Find the hover equilibrium by substitution}
An equilibrium has zero time derivatives. Choose $v_0=0$, $\omega_0=0$,
$\phi_0=\theta_0=0$ and any fixed heading $\psi_0$. Substitution in the
vertical equation gives $0=T_0-mg$, so $T_0=mg$. Zero torque with a symmetric
allocation gives $T_{i0}=mg/4$. If normalized desired thrust satisfies
$T_i=T_{\max}u_i$, then
\begin{equation}
u_{i0}=c_0=\frac{mg}{4T_{\max}}.
\end{equation}
For $m=1.2$ kg and $T_{\max}=6$ N, each motor must supply $2.941995$ N and
$c_0=0.4903325$. This is a model calculation, not an identified ESC duty cycle.
The thrust-to-weight ratio of this assumed model is $24/11.76798\simeq2.0394$.
It is not evidence that the real 1400KV/9047 combination has that reserve.

\section{Linearize by explicitly retaining first-order terms}
Define perturbations $\delta T=T-T_0$, $\delta\phi=\phi-\phi_0$ and so on.
The first-order Taylor expansion of $\sin\epsilon$ is $\epsilon$; that of
$\cos\epsilon$ is $1$. Products such as $\delta T\delta\theta$ contain two
small perturbations and are second order. At heading zero,
\begin{align}
(T_0+\delta T)\sin\delta\theta\cos\delta\phi
 &\simeq T_0\delta\theta,\\
-(T_0+\delta T)\sin\delta\phi&\simeq-T_0\delta\phi,\\
(T_0+\delta T)\cos\delta\theta\cos\delta\phi-mg&\simeq\delta T.
\end{align}
Using $T_0/m=g$, define $\gamma=d_v/m=0.1$ s$^{-1}$ and obtain
\begin{align}
\delta\ddot x+\gamma\,\delta\dot x&=g\,\delta\theta+f_{wx}/m,\\
\delta\ddot y+\gamma\,\delta\dot y&=-g\,\delta\phi+f_{wy}/m,\\
\delta\ddot h+\gamma\,\delta\dot h&=\delta T/m+f_{wz}/m.
\end{align}
Here $\gamma$ is the damping coefficient and $\delta x$ denotes a perturbation.
Below we omit the perturbation prefix to keep the equations legible. At a nonzero heading, rotate the two
horizontal acceleration equations by the heading angle. This is a local
linearization, not a claim that Euler coordinates work globally.

The body-rate products $qr,rp,pq$ vanish to first order about zero rates.
Near level attitude $\dot\phi\simeq p$, $\dot\theta\simeq q$, $\dot\psi\simeq r$.
Thus $J_x\ddot\phi=\tau_x$, $J_y\ddot\theta=\tau_y$ and
$J_z\dot r=\tau_z$. There is intentionally no yaw-angle feedback requirement.

\section{Build each open-loop transfer function from the differential equation}
A transfer function describes a linear input/output relation with zero initial
perturbations. For $\dot y=f$, the unilateral Laplace transform is
$sY(s)-y(0)=F(s)$. For $\ddot y$, it is
$s^2Y(s)-sy(0)-\dot y(0)$. Setting the initial perturbations to zero removes
those initial-condition terms; it does not mean the actual aircraft has no
initial state. Initial-condition responses must be added separately.

From $J_x\ddot\phi=\tau_x$, first transform to
$J_xs^2\Phi=\mathcal{T}_x$, then divide both sides by $J_xs^2\mathcal{T}_x$:
\begin{equation}
\frac{\Phi(s)}{\mathcal{T}_x(s)}=\frac{1}{J_xs^2},\quad
\frac{P(s)}{\mathcal{T}_x(s)}=\frac{1}{J_xs}.
\end{equation}
Use the allocation $\tau_x=4aT_{\max}u_\phi$ and the motor lag
$\tau_m\dot u_{\rm act}+u_{\rm act}=u_{\rm cmd}$.
Transforming the latter gives
$(\tau_ms+1)U_{\rm act}=U_{\rm cmd}$. Therefore the cascade is
\begin{equation}
G_{p}(s)=\frac{P_d(s)}{U_\phi(s)}=
\frac{b_x}{s(\tau_ms+1)},\qquad
G_{\phi}(s)=\frac{\Phi_d(s)}{U_\phi(s)}=
\frac{b_x}{s^2(\tau_ms+1)},
\end{equation}
where $b_x=(4aT_{\max}/J_x)(180/\pi)$ because the controller uses degrees.
Pitch has $b_y$ instead; yaw rate has $b_\psi=4\kappa T_{\max}(180/\pi)/J_z$.
The rate plant is second order with lag and the angle plant third order.
Without actuator lag these reduce to first and second order respectively.

For altitude, $\ddot h+\gamma\dot h=b_hu_{\rm act}$ with
$b_h=4T_{\max}/m=20$ m/s$^2$/command. Transforming each derivative gives
\begin{equation}
G_h(s)=\frac{H(s)}{U_h(s)}=
\frac{b_h}{s(s+\gamma)(\tau_ms+1)}.
\end{equation}
This retains drag and actuator lag. The often-used double integrator $b_h/s^2$
is obtained only after explicitly neglecting both. Horizontal position from
small angle is $X/\Theta=g/[s(s+\gamma)]$ and
$Y/\Phi=-g/[s(s+\gamma)]$ in radians. Degree-valued angle inputs require an
additional factor $\pi/180$.

\section{Relate normalized thrust to a possible rotor-speed model}
A better identified actuator could use $T=k_T\Omega^2$ and
$Q=k_Q\Omega^2$. Around nonzero hover speed $\Omega_0$, expand
\begin{equation}
k_T(\Omega_0+\Delta\Omega)^2
=k_T\Omega_0^2+2k_T\Omega_0\Delta\Omega+k_T(\Delta\Omega)^2.
\end{equation}
Discarding the quadratic perturbation gives
$\Delta T=2k_T\Omega_0\Delta\Omega$. A first-order identified command-to-speed
gain $k_\Omega/(\tau_\Omega s+1)$ then yields
$\Delta T/\Delta u=2k_T\Omega_0k_\Omega/(\tau_\Omega s+1)$.
The current normalized-thrust model replaces that unknown product with
$T_{\max}$. KV alone supplies none of $k_T,k_Q,k_\Omega,\tau_\Omega$ under
load. A higher-order model is useful only when its extra parameters are measured.

\chapter{Closed-loop transfer functions and stability calculations}
\section{Rate PI: derive the denominator without skipping the feedback algebra}
Let $G=b/[s(\tau_ms+1)]$, $C=K_p+K_i/s$ and $e=r_c-r$.
Then $r=GCe=GC(r_c-r)$. Move the feedback term to the left:
$(1+GC)r=GC r_c$. Divide by $1+GC$ and clear denominators:
\begin{equation}
T_r(s)=\frac{r}{r_c}=
\frac{b(K_ps+K_i)}{\tau_ms^3+s^2+bK_ps+bK_i}.
\end{equation}
The numerator zero at $s=-K_i/K_p$ changes the step response; denominator
poles alone do not completely describe overshoot. With $\tau_m=0$, matching
$s^2+2\zeta\omega_ns+\omega_n^2$ gives $K_p=2\zeta\omega_n/b$ and
$K_i=\omega_n^2/b$. This is an initial design followed by checking the cubic,
not an exact second-order description of the complete plant.

\section{Roll/pitch angle cascade: retain the inner loop dynamics}
Let $k_a=3$ s$^{-1}$ and $p_c=k_a(\phi_c-\phi)$. Since $p=s\phi$,
\begin{align}
\phi&=\frac{T_r}{s}k_a(\phi_c-\phi),\\
\left(1+\frac{k_aT_r}{s}\right)\phi&=\frac{k_aT_r}{s}\phi_c,\\
T_\phi(s)&=\frac{k_aT_r}{s+k_aT_r}.
\end{align}
Substitute the rate-loop numerator and denominator and collect powers:
\begin{equation}
\boxed{T_\phi(s)=\frac{bk_a(K_ps+K_i)}{
\tau_ms^4+s^3+bK_ps^2+b(K_i+k_aK_p)s+bk_aK_i}.}
\end{equation}
Pitch is identical only under the assumed $J_x=J_y$. This fourth-order
expression is more complete than replacing the inner rate loop by unity.
The ideal outer approximation $k_a/(s+k_a)$ is justified only well below the
inner closed-loop bandwidth, away from saturation and estimator error.

\section{Altitude PID and the distinction between height and velocity requests}
At fixed equilibrium, $e_h=h_c-h$, $\dot I_h=e_h$ and
$u=K_pe_h+K_iI_h+K_d(v_c-v)$, where the internally integrated target satisfies
$\dot h_c=v_c$. For differentiable consistent references,
$V_c=sH_c$, $V=sH$ and $I_h=(H_c-H)/s$, so
\begin{equation}
U=(K_p+K_i/s+K_ds)(H_c-H).
\end{equation}
Combine this with $G_h$ above and multiply the characteristic equation by
$s^2(s+\gamma)(\tau_ms+1)$:
\begin{align}
D_h(s)&=s^2(s+\gamma)(\tau_ms+1)
+b_h(K_ds^2+K_ps+K_i)\\
&=\tau_ms^4+(1+\gamma\tau_m)s^3
+(\gamma+b_hK_d)s^2+b_hK_ps+b_hK_i,\\
\boxed{T_h(s)}&=\boxed{\frac{b_h(K_ds^2+K_ps+K_i)}{D_h(s)}}.
\end{align}
The transfer from commanded vertical velocity to height is $T_h(s)/s$;
the transfer to vertical velocity is $T_h(s)$, with zero initial perturbations.
A discontinuous height step accompanied by an impulsive derivative reference
is not the actual stick interface. The mission generates bounded $v_c$; the
height target integrates it. If one instead steps $h_c$ while forcing $v_c=0$,
the numerator is $b_h(K_ps+K_i)$: the denominator is the same, but the reference
response differs. This prevents applying a textbook PID step response to the
wrong command interface.

When lag is neglected but drag retained, matching
$(s+p_h)(s^2+2\zeta_h\omega_hs+\omega_h^2)$ gives
\begin{align}
K_d&=(2\zeta_h\omega_h+p_h-\gamma)/b_h,\\
K_p&=(\omega_h^2+2p_h\zeta_h\omega_h)/b_h,\\
K_i&=p_h\omega_h^2/b_h.
\end{align}
The existing baseline was designed with $\gamma=0$, giving $K_d=0.175$.
Including the assumed drag gives $K_d=0.170$. The selected $\lambda=1.3$
candidate similarly gives $0.2275$ versus $0.2225$. These differences are shown
as model alternatives; gains are not silently changed without a new route test.

\section{What second-order parameters actually mean}
For a complex pair $s_{1,2}=-\sigma\pm j\omega_d$ with $\sigma>0$,
$\omega_n=\sqrt{\sigma^2+\omega_d^2}$ and $\zeta=\sigma/\omega_n$.
For the ideal normalized second-order low-pass without zeros,
\begin{equation}
M_p=e^{-\pi\zeta/\sqrt{1-\zeta^2}},\quad
t_p=\pi/\omega_d,\quad t_{s,2\%}\simeq4/(\zeta\omega_n).
\end{equation}
These follow from the decaying sinusoidal homogeneous response and its extrema;
the settling approximation bounds its exponential envelope approximately.
For example the ideal rate design $\zeta=0.9$, $\omega_n=12$ has
$\omega_d=5.2307$ rad/s and envelope settling about $0.3704$ s. Its ideal
no-zero overshoot is about $0.152\%$. The actual PI closed loop has a zero,
and actuator lag adds another pole. Therefore these figures describe the
chosen target polynomial, not a guaranteed measured aircraft step response.
The companion numerical analysis computes actual step responses for the
complete local transfer functions, including their zeros.

\section{Routh--Hurwitz: construct and evaluate the stability conditions}
For $a_3s^3+a_2s^2+a_1s+a_0$ with $a_3>0$, the Routh first column is
\begin{equation}
a_3,\quad a_2,\quad\frac{a_2a_1-a_3a_0}{a_2},\quad a_0.
\end{equation}
All entries must be strictly positive for strict left-half-plane roots, assuming
no degenerate zero rows. Applying this to the rate loop gives positive gains
and $K_p>\tau_mK_i$. Substituting the pole-matched PI gives
$\omega_n<2\zeta/\tau_m$. For $\tau_m=0.04$ s the bound is 45 rad/s;
the selected rate frequency 15.6 rad/s meets this local cubic condition.

For $a_4s^4+a_3s^3+a_2s^2+a_1s+a_0$, construct
\begin{equation}
\begin{array}{c|ccc}
s^4&a_4&a_2&a_0\\s^3&a_3&a_1&0\\
s^2&\Delta_2/a_3&a_0&0\\
s^1&\Delta_3/\Delta_2&0&0\\s^0&a_0&0&0
\end{array},
\end{equation}
where $\Delta_2=a_3a_2-a_4a_1$ and
$\Delta_3=a_1\Delta_2-a_3^2a_0$. With $a_4,a_3,a_0>0$, strict stability
requires $\Delta_2>0$ and $\Delta_3>0$. Positive polynomial coefficients alone
are insufficient. Substitution for the full angle cascade gives
\begin{align}
\Delta_2&=bK_p-\tau_mb(K_i+k_aK_p),\\
\Delta_3&=b(K_i+k_aK_p)\Delta_2-bk_aK_i.
\end{align}
For altitude use $a_3=1+\gamma\tau_m$, $a_2=\gamma+b_hK_d$,
$a_1=b_hK_p$ and $a_0=b_hK_i$. The executable analysis reports each
determinant, every pole and the actual step metrics for both gain sets.
These strict tests apply to the local unsaturated models. Ground contact,
mode transitions and a changing camera map are not a single LTI system.

\section{A continuous parameter-interval proof for the conservative candidate}
The grid finds counterexamples. For $\lambda=0.7$ we can additionally prove
local Routh inequalities over a continuous box. Assume
$b\in[0.65b_0,1.35b_0]$, $\tau\in[0.02,0.09]$ s, and for height
$\gamma\in[0.08,0.13]$ s$^{-1}$. These are assumed ranges, not measured
confidence intervals for the real aircraft.

For the angle cascade define $u=K_i+k_aK_p>0$ and
$A(\tau)=K_p-\tau u$. Since $u>0$, $A$ decreases with $\tau$, so
$A\geq A_{\min}=K_p-\tau_{\max}u$. The determinants become
\begin{equation}
\Delta_2=bA,\qquad\Delta_3=b[buA-k_aK_i].
\end{equation}
If $A_{\min}>0$, multiplication by positive $b,u$ preserves inequalities.
It is therefore sufficient to check
\begin{equation}
A_{\min}>0,\qquad b_{\min}uA_{\min}-k_aK_i>0.
\end{equation}
The numerical lower bounds are $0.000257095$ and $0.00776085$ approximately.
All coefficients and Routh first-column entries are consequently positive
throughout that continuous box for this decoupled angle model. Yaw requires
$K_p-\tau_{\max}K_i>0$, with lower bound $0.00462303$.

For height, expand the determinant without skipping terms:
\begin{align}
\Delta_2&=(1+\gamma\tau)(\gamma+bK_d)-\tau bK_p\\
&=\gamma+\gamma^2\tau+bK_d+b\gamma\tau K_d-b\tau K_p.
\end{align}
Since the discarded $\gamma^2\tau$ and $b\gamma\tau K_d$ terms are nonnegative,
and $K_d-\tau_{\max}K_p>0$ for this candidate,
\begin{equation}
\Delta_2\geq\gamma_{\min}+b_{\min}(K_d-\tau_{\max}K_p)
 =D_{2,\rm lb}>0.
\end{equation}
The next determinant obeys
\begin{align}
\Delta_3&=b\left[K_p\Delta_2-K_i(1+\gamma\tau)^2\right]\\
&\geq b\left[K_pD_{2,\rm lb}-K_i(1+\gamma_{\max}\tau_{\max})^2\right].
\end{align}
Here $D_{2,\rm lb}\simeq1.5460874$ and the final bracket is approximately
$0.1354504>0$. This proves strict local height-loop stability over the specified
continuous box. It excludes saturation, nonlinear coupling, estimator dynamics,
contacts, delays and companion guidance. Those exclusions are part of the theorem.

The fast candidate also has an analytic counterexample. For $\lambda=1.3$,
$K_i/K_p=8.6667$ s$^{-1}$. The necessary angle condition $A>0$ implies
$k_a<1/\tau-K_i/K_p$. At $\tau=0.09$ s the right side is
$2.4444$ s$^{-1}$, below the chosen $k_a=3$. Thus this Routh condition fails
regardless of positive plant gain at that lag. Nominal success at 0.04 s
does not establish stability at 0.09 s.

\section{The horizontal position loop is sixth order}
The companion generates $a_c=k_x(x_c-x)-k_v\dot x$ and $\theta_c=a_c/g$
in radians. Use $X/\Theta=g/[s(s+\gamma)]$ and the complete angle cascade
$T_\theta=N_\theta/D_\theta$. Substitution gives
\begin{align}
X&=\frac{T_\theta}{s(s+\gamma)}[k_xX_c-(k_x+k_vs)X],\\
\big[s(s+\gamma)D_\theta+(k_vs+k_x)N_\theta\big]X&=k_xN_\theta X_c,\\
\boxed{\frac{X}{X_c}}&=\boxed{\frac{k_xN_\theta}
 {s(s+\gamma)D_\theta+(k_vs+k_x)N_\theta}}.
\end{align}
This denominator is sixth order. Navigation uses $k_x=1.8$ s$^{-2}$,
$k_v=1.8$ s$^{-1}$; approach/landing uses $k_x=1.2$ and the same $k_v$.
The lateral-roll sign cancels with its guidance-conversion sign, giving the
same local scalar polynomial under the symmetric model. This is upstream
guidance; there are still exactly six flight-controller PID objects.

For $D_x=d_6s^6+d_5s^5+\cdots+d_0$, the specific Hurwitz array is
\begin{equation}
\mathcal H=\begin{bmatrix}
d_5&d_3&d_1&0&0&0\\d_6&d_4&d_2&d_0&0&0\\
0&d_5&d_3&d_1&0&0\\0&d_6&d_4&d_2&d_0&0\\
0&0&d_5&d_3&d_1&0\\0&0&d_6&d_4&d_2&d_0
\end{bmatrix}.
\end{equation}
For $d_6>0$ all six leading principal determinants must be positive. The
executable analysis reports every determinant and pole for each candidate.
Speed limits, changing waypoints and camera-map updates remain nonlinear
effects beyond this local calculation.

\section{Discrete-time stability at the real control interval}
Keep the physical state $x_p=[\phi,p,u_{\rm act}]^{\mathsf T}$ and integral
state $I$. During one held command, obtain the exact linear plant update
$x_{p,k+1}=A_dx_{p,k}+B_du_k$ by solving the linear ODE over 2 ms.
The software uses the current error in its candidate integral:
\begin{align}
e_k&=k_a(\phi_{c,k}-\phi_k)-p_k,\\
I_{k+1}&=I_k+\Delta e_k,\\
u_k&=K_pe_k+K_iI_{k+1}=(K_p+K_i\Delta)e_k+K_iI_k.
\end{align}
Let $C_e=[-k_a,-1,0]$. At zero reference the combined update is therefore
\begin{equation}
\begin{bmatrix}x_{p,k+1}\\I_{k+1}\end{bmatrix}=
\begin{bmatrix}
A_d+B_d(K_p+K_i\Delta)C_e&B_dK_i\\
\Delta C_e&1
\end{bmatrix}
\begin{bmatrix}x_{p,k}\\I_k\end{bmatrix}.
\end{equation}
This explicitly includes the integration ordering used in C. Strict sampled
stability requires every eigenvalue $z_i$ to satisfy $|z_i|<1$. Yaw uses the
same construction without the angle state; altitude uses $C_e=[-1,0,0]$
and velocity feedback $-K_dv$. Estimator dynamics, computation delay and
allocation limits are additional effects not present in this particular test.
The numerical analysis also repeats it with one control sample of input delay
and reports the resulting spectral radius rather than asserting timing safety.

\section{Include the 20 Hz companion with the 500 Hz flight loop}
For one horizontal axis use physical state
$x_p=[\theta_d,p_d,u_{\rm act},x,v]^{\mathsf T}$. Its hover plant is
\begin{equation}
\dot x_p=\begin{bmatrix}
0&1&0&0&0\\0&0&b_x&0&0\\0&0&-1/\tau_m&0&0\\
0&0&0&0&1\\g\pi/180&0&0&0&-\gamma
\end{bmatrix}x_p+
\begin{bmatrix}0\\0\\1/\tau_m\\0\\0\end{bmatrix}u.
\end{equation}
The $g\pi/180$ entry explicitly converts the degree-valued angle into the
radian-valued small-angle force law. Obtain $A_d,B_d$ for one held 2 ms
command, append the PI integral, and apply the same current-error integral
ordering as above. With $k=K_p+K_i\Delta$, $C_e=[-k_a,-1,0,0,0]$,
\begin{equation}
F=\begin{bmatrix}A_d+B_dkC_e&B_dK_i\\\Delta C_e&1\end{bmatrix},\qquad
G=\begin{bmatrix}B_dkk_a\\\Delta k_a\end{bmatrix}.
\end{equation}
Thus $x_{k+1}=Fx_k+G\theta_c$ for the six-state augmented system.
The companion holds $\theta_c=C_gx_k$ for $N=25$ flight updates, where
\begin{equation}
C_g=-\frac{180}{\pi g}[0,0,0,k_x,k_v,0].
\end{equation}
Substitute the recurrence a second time:
$x_{k+2}=F^2x_k+(F+I)G\theta_c$. A third substitution gives
$F^3x_k+(F^2+F+I)G\theta_c$. Continuing the same physical recurrence for
25 held steps gives
\begin{equation}
x_{k+N}=F^Nx_k+\left(\sum_{j=0}^{N-1}F^jG\right)C_gx_k=Mx_k.
\end{equation}
The numerical analysis checks every eigenvalue of $M$. With one complete
guidance-period delay, append the old held command and check
\begin{equation}
M_{\rm delayed}=\begin{bmatrix}
F^N&\sum_{j=0}^{N-1}F^jG\\C_g&0
\end{bmatrix}.
\end{equation}
This is an explicit 500 Hz/20 Hz sampled-system calculation. It still assumes
exact state feedback, linear hover dynamics and scheduled updates. It is not
a proof for the Kalman estimator, delayed perception, saturation, or Windows
meeting every real-time deadline.

\section{Frequency-domain margins and uncertainty}
For the rate loop, open-loop $L_r=C G=b(K_ps+K_i)/[s^2(\tau_ms+1)]$.
At the gain crossover $|L(j\omega_c)|=1$, phase margin is
$180^\circ+\arg L(j\omega_c)$. At a $-180^\circ$ phase crossover,
gain margin is $1/|L(j\omega_\pi)|$. These are frequency-domain local robustness
measures, not obstacle-clearance margins. A pure delay adds phase
$-\omega T_d$; the rough additional-delay allowance from a positive phase
margin is $T_d\simeq\mathrm{PM}_{\rm rad}/\omega_c$ for the crossing being
considered. Multiple crossovers and unstable open-loop cases need care.

Plant uncertainty enters directly: $b_x\propto aT_{\max}/J_x$ and
$b_h\propto T_{\max}/m$. The executable analysis varies these gains and
$\tau_m$ over a declared finite grid and counts stable/unstable cases for the
complete characteristic polynomials. A grid with no failures is evidence
only at those points; it is not a continuous robust-stability proof. Battery
sag, thrust asymmetry, flexible arms and estimator bias remain unmodelled
physical risks until measured. Time-domain runs then test saturation, takeoff,
landing and perception-driven commands beyond the linear analysis.

\chapter{Companion missions, boost and actuator authority}
\section{Keep one flight-control mode; move mission sequencing upstream}
The flight core still receives roll/pitch angles, yaw rate and centered
vertical-speed throttle. The companion selects mission phases: ground,
spool-up in simulation, takeoff, hover, navigate, approach/hold, descent and
landed. A mission phase does not introduce a new flight PID. The companion
converts a requested height into bounded vertical speed; the existing altitude
PID integrates that speed into its internal target. Start and finish points
are explicit local coordinates. The start must agree with the measured launch
position; selecting it must never teleport the aircraft state.

Ground contact supplies an upward normal reaction that balances weight before
lift-off. At $T_z>mg$ that unilateral reaction vanishes and upward acceleration
begins. Landing is the reverse transition with low descent speed and contact
confirmation; simply forcing a trajectory through $h=0$ is not a contact model.
The 3D simulator handles contact; real landed detection remains a separate
sensor/state-machine problem. A simulation-only motor gate cannot be presented
as an implemented physical arming system.

\section{What arming and RT boost mean}
Arming is a deliberate permission to produce physical motor outputs. It is not
the same as compiling a control function or connecting a battery. Physical
arming stays disabled as requested. Takeoff/landing button gestures are mission
requests; a future authorized safety supervisor must accept them before motors
can be enabled. The simulated actuator gate is explicitly isolated from hardware.

RT is already present in the radio frame as a normalized trigger. Its proposed
boost increases the maximum commanded vertical speed from 2 to 3 m/s and the
nominal collective ceiling from 0.65 to 0.80 as RT moves from zero to one.
The internal boost allowance rises over at least 0.25 s and decays over at least
0.10 s for a full-scale change. RT alone does not command a climb at centered
vertical stick. The companion simulation inhibits boost during approach/landing.
Roll/pitch leveling and yaw-rate control remain active. It is not acrobatic
rate mode and not an unconditional full-throttle override. The boost ceiling
has less correction reserve than normal mode; maximum climb requires measured
thrust/energy reserve and later hardware tuning.

\section{Prove the motor-headroom allocation interval}
Let desired motor values before clipping be $u_i=c+d_i$, where $d_i$ are the
four differential roll/pitch/yaw corrections. For each motor to satisfy
$0\leq u_i\leq1$, rearrange to $-d_i\leq c\leq1-d_i$.
Intersect all four intervals:
\begin{equation}
-\min_i d_i\leq c\leq1-\max_i d_i.
\end{equation}
This interval exists exactly when $\max d_i-\min d_i\leq1$. If the requested
differential span exceeds one, uniformly scale all $d_i$ by
$\alpha=1/(\max d_i-\min d_i)$; otherwise $\alpha=1$. Clamp collective to
the resulting feasible interval and calculate $u_i=c+\alpha d_i$.
This preserves differential-command ratios, unlike clipping each motor
independently. It cannot create missing thrust or guarantee full simultaneous
torque authority. A future allocator should feed allocation residuals back into
anti-windup; the current PID anti-windup covers each loop's own output limit,
not the complete coupled allocation constraint.

\chapter{Exactly six feedback controllers}
\section{Signals and units}
The six objects in the C file are roll angle, roll rate, pitch angle, pitch
rate, yaw rate and altitude. A PID object may have a zero integral or
derivative coefficient. Calling all six ``PIDs'' describes the general
structure; the chosen outer loops are P and the chosen angular-rate loops
are PI. No seventh vertical-rate PID is hidden in the altitude loop.

Angular gains below use \textbf{degrees} and degrees per second because that
is the C interface. Physics uses radians. If an angular measurement is
changed from degrees to radians, multiply corresponding numeric gains by
$180/\pi$ to preserve the same normalized control output. Outer angle-to-rate
P gains remain unchanged when both angle and rate units are converted together.

For a generic PID,
\begin{align}
e_k&=r_k-y_k,\\
I_k^{\rm cand}&=\operatorname{clip}(I_{k-1}+e_k\Delta,-I_{\max},I_{\max}),\\
u_k^{\rm raw}&=K_pe_k+K_i I_k^{\rm cand}-K_d\dot y_k,\\
u_k&=\operatorname{clip}(u_k^{\rm raw},-u_{\max},u_{\max}).
\end{align}
The derivative on measurement avoids setpoint derivative kick. All selected
angular $K_d$ values are zero. The altitude controller instead uses the
available derivative reference $v_c$: $K_d(v_c-\hat v_z)$.
The integrator candidate is accepted if unsaturated, or if the error drives
an existing saturation back toward the admissible interval. Integrator and
output limits are independent physical quantities.

\section{Roll and pitch rate PI derivation}
Near hover neglect cross coupling and lag for the initial design. With
angular rate measured in deg/s,
\begin{equation}
\dot p_d=b_xu_\phi,\quad b_x=\frac{4aT_{\max}}{J_x}\frac{180}{\pi},\qquad
\dot q_d=b_yu_\theta,\quad b_y=\frac{4aT_{\max}}{J_y}\frac{180}{\pi}.
\end{equation}
The plant transfer function from normalized torque command to rate is $b/s$.
The PI controller is $C(s)=K_p+K_i/s=(K_ps+K_i)/s$. Unity feedback gives
\begin{align}
1+C(s)G(s)&=0,\\
1+\frac{b(K_ps+K_i)}{s^2}&=0,\\
s^2+bK_ps+bK_i&=0.
\end{align}
Choose a desired second-order polynomial $s^2+2\zeta\omega_n s+\omega_n^2$.
Matching like powers of $s$ proves
\begin{equation}\boxed{K_p=\frac{2\zeta\omega_n}{b},\qquad K_i=\frac{\omega_n^2}{b},\qquad K_d=0.}\end{equation}
For the baseline $\zeta=0.9$, $\omega_n=12$ rad/s. The roots have real part
$-\zeta\omega_n<0$. A state-space proof is also possible: with error
$e=r-p_d$ and integral $I$, for constant $r$,
\begin{equation}
\frac{\dd}{\dd t}\begin{bmatrix}e\\I\end{bmatrix}
=\begin{bmatrix}-bK_p&-bK_i\\1&0\end{bmatrix}
\begin{bmatrix}e\\I\end{bmatrix}.
\end{equation}
Its determinant $\det(sI-A)=s^2+bK_ps+bK_i$ gives the same condition.

Including the actuator lag changes the characteristic polynomial to
\begin{equation}
\tau_m s^3+s^2+bK_ps+bK_i.
\end{equation}
For positive coefficients, the cubic Routh condition additionally requires
$bK_p>\tau_m bK_i$, or $K_p>\tau_mK_i$.
With the pole-matched gains this is $\omega_n<2\zeta/\tau_m=45$ rad/s.
Both the baseline 12 and selected 15.6 rad/s designs meet this continuous
linear condition. It is not a proof for saturation, delays, estimators or the
full nonlinear aircraft.

\section{Outer angle loops and yaw}
The outer commands are
\begin{equation}
p_c=3(\phi_c-\hat\phi),\qquad q_c=3(\theta_c-\hat\theta)
\quad[\mathrm{deg/s}],
\end{equation}
clamped to $\pm180$ deg/s. Under the ideal approximation $p_d=p_c$,
$\dot e_\phi=-3e_\phi$, so $e_\phi(t)=e_\phi(0)e^{-3t}$.
The time constant is $1/3$ s and the ideal two-percent settling time is
$-\ln(0.02)/3\simeq1.304$ s. Real cascaded dynamics are slower/different.
The inner angular bandwidth is intentionally higher than this outer loop.

Yaw has \emph{only} rate feedback in the flight core:
\begin{equation}
\dot r_d=b_\psi u_\psi,\quad
b_\psi=\frac{4\kappa T_{\max}}{J_z}\frac{180}{\pi},\quad
K_{p\psi}=\frac{2(0.9)\omega_\psi}{b_\psi},\quad
K_{i\psi}=\frac{\omega_\psi^2}{b_\psi}.
\end{equation}
Baseline $\omega_\psi=4$ rad/s. There is no yaw-angle PID in the flight
controller. The companion optionally creates a bounded yaw-rate request from
its local-heading error; that is a navigation policy, not an extra flight PID.

\section{One altitude PID with vertical-speed stick command}
For normalized centered throttle $s_t\in[0,1]$,
\begin{equation}
v_c=4(s_t-0.5)\quad\mathrm{m/s},\qquad
h_{c,k+1}=h_{c,k}+v_c\Delta.
\end{equation}
At centre stick, $v_c=0$ and the target height is held. The target is initialized
to the estimated height and limited to within 2 m of it to avoid unlimited
reference accumulation. There is one altitude control law:
\begin{align}
e_h&=h_c-\hat h,\quad\dot I_h=e_h,\\
\delta c&=K_{ph}e_h+K_{ih}I_h+K_{dh}(v_c-\hat v_z),\\
c&=\operatorname{clip}\left(\frac{c_0+\delta c}
 {\max(0.5,\cos\hat\phi\cos\hat\theta)},0,1\right).
\end{align}
The tilt factor approximately preserves vertical thrust. The factor limit
prevents numerical blow-up; it is not valid recovery control at large tilt.
The nominal hover command is
\begin{equation}c_0=\frac{mg}{4T_{\max}}=0.4903325.\end{equation}
The board scaffold still has a 0.50 placeholder; SITL explicitly supplies
the model-based value. A real craft requires a measured hover thrust curve.

At fixed target, near hover, let $b_h=4T_{\max}/m=20$ m/s$^2$ per command.
Then $\ddot h=b_h\delta c$. Substitution of PID feedback gives
\begin{equation}
s^3+b_hK_{dh}s^2+b_hK_{ph}s+b_hK_{ih}=0.
\end{equation}
Choose $(s+p_h)(s^2+2\zeta_h\omega_hs+\omega_h^2)$, multiply explicitly:
\begin{align}
&s^3+2\zeta_h\omega_hs^2+\omega_h^2s
+p_hs^2+2p_h\zeta_h\omega_hs+p_h\omega_h^2\\
&=s^3+(2\zeta_h\omega_h+p_h)s^2
+(\omega_h^2+2p_h\zeta_h\omega_h)s+p_h\omega_h^2.
\end{align}
Coefficient matching yields
\begin{equation}\boxed{\begin{aligned}
K_{dh}&=(2\zeta_h\omega_h+p_h)/b_h,\\
K_{ph}&=(\omega_h^2+2p_h\zeta_h\omega_h)/b_h,\\
K_{ih}&=p_h\omega_h^2/b_h.
\end{aligned}}\end{equation}
Baseline $\omega_h=1.5$, $\zeta_h=0.9$, $p_h=0.8$ gives
$(K_p,K_i,K_d)=(0.2205,0.09,0.175)$. Units are respectively
command/m, command/(m s), and command s/m. The integral-state bound is
2 m s and the altitude-correction bound is 0.35 command.
Actuator lag increases the order:
\begin{equation}
\tau_m s^4+s^3+b_hK_d s^2+b_hK_p s+b_hK_i.
\end{equation}
For positive coefficients the quartic Routh conditions reduce to
\begin{equation}
b_hK_d>\tau_m b_hK_p,\qquad
(b_hK_d)(b_hK_p)>\tau_m(b_hK_p)^2+b_hK_i.
\end{equation}
These hold for the baseline and selected simulation gains. Discrete feedback
and filtering still require the actual closed-loop tests.

\section{A reproducible gain family, not hidden hand tuning}
Three candidates use scale $\lambda\in\{0.7,1,1.3\}$ with
$\omega_n=12\lambda$, $\omega_\psi=4\lambda$,
$\omega_h=1.5\lambda$, $p_h=0.8\lambda$; outer P stays 3.
Thus rate $K_p$ scales as $\lambda$, rate $K_i$ as $\lambda^2$,
altitude $K_p$ as $\lambda^2$, altitude $K_i$ as $\lambda^3$,
and altitude $K_d$ as $\lambda$. Full numeric tables and the declared selection
score appear in the results chapter. The selected family is the best of these
three under that score, not a globally optimal tune.

\chapter{Kalman filtering: proof and implemented equations}
\section{Derivation of the gain}
Let $\hat x^-$ be the prior estimate, $e^-=x-\hat x^-$, with covariance
$P^-$. A measurement is $z=Hx+v$ with independent noise covariance $R$.
The innovation is $\nu=z-H\hat x^-=He^-+v$. Hence
\begin{equation}S=\operatorname{Cov}(\nu)=HP^-H^{\mathsf T}+R.\end{equation}
Use $\hat x^+=\hat x^-+K\nu$. Then
\begin{equation}
e^+=(I-KH)e^- - Kv.
\end{equation}
Taking its covariance, retaining both independent-noise contributions, gives
the Joseph form
\begin{equation}\label{eq:joseph}
P^+=(I-KH)P^-(I-KH)^{\mathsf T}+KRK^{\mathsf T}.
\end{equation}
Expanding this expression gives
$P^+=P^- -KHP^- -P^-H^{\mathsf T}K^{\mathsf T}+KSK^{\mathsf T}$.
The expected squared error is $\E[e^{+\mathsf T}e^+]=\operatorname{tr}(P^+)$.
Complete the square using $K_*=P^-H^{\mathsf T}S^{-1}$:
\begin{equation}
P^+=P^- -P^-H^{\mathsf T}S^{-1}HP^-
 +(K-K_*)S(K-K_*)^{\mathsf T}.
\end{equation}
The last term is positive semidefinite when $S$ is positive definite. Its
trace is minimized by $K=K_*$, proving the Kalman gain. Gaussianity is not
required for this best linear update; Gaussian priors/noises additionally
make it the conditional-mean Bayesian update and justify normal intervals.

\section{Scalar one-angle filter}
A strictly scalar attitude model (calibrated gyro bias assumed known) is
\begin{align}
\hat\theta^-_k&=\hat\theta^+_{k-1}+\Delta(\omega_{g,k}-b),\\
P^-_k&=P^+_{k-1}+\Delta^2\sigma_g^2,\\
K_k&=\frac{P^-_k}{P^-_k+R_\theta},\\
\hat\theta^+_k&=\hat\theta^-_k+K_k(\theta_{a,k}-\hat\theta^-_k),\\
P^+_k&=(1-K_k)^2P^-_k+K_k^2R_\theta.
\end{align}
This is the literal one-state interpretation of ``1D Kalman''. $K$ is not
an arbitrary fixed mixing percentage: it changes with prior uncertainty,
sample interval and measurement noise.

\section{Actual MPU filter: one axis, two states}
The C implementation estimates \emph{one angle at a time}, independently
for roll and pitch, but also estimates gyro bias. Its state dimension is
therefore two, even though it is often informally called a ``1D attitude
filter''. This document names that distinction explicitly.
\begin{equation}
x=\begin{bmatrix}\theta\\b\end{bmatrix},\quad
F=\begin{bmatrix}1&-\Delta\\0&1\end{bmatrix},\quad
B=\begin{bmatrix}\Delta\\0\end{bmatrix},\quad H=\begin{bmatrix}1&0\end{bmatrix}.
\end{equation}
The bias is subtracted from the gyro, so the upper-right entry is negative.
With independent per-sample gyro variance $\sigma_g^2$ and continuous bias
random-walk spectral density $q_b$,
\begin{equation}
Q=\begin{bmatrix}
\sigma_g^2\Delta^2+q_b\Delta^3/3&-q_b\Delta^2/2\\
-q_b\Delta^2/2&q_b\Delta
\end{bmatrix}.
\end{equation}
The bias contribution follows by integrating
$q_b[-\tau\ 1]^{\mathsf T}[-\tau\ 1]$ from zero to $\Delta$. Expanding
$FPF^{\mathsf T}$ entry by entry gives
\begin{align}
P^-_{00}&=P_{00}-\Delta(P_{01}+P_{10})+\Delta^2P_{11}
 +\sigma_g^2\Delta^2+q_b\Delta^3/3,\\
P^-_{01}&=P_{01}-\Delta P_{11}-q_b\Delta^2/2,\\
P^-_{10}&=P_{10}-\Delta P_{11}-q_b\Delta^2/2,\\
P^-_{11}&=P_{11}+q_b\Delta.
\end{align}
Since $H=[1,0]$, the matrix inverse in the gain formula is just a scalar
division:
\begin{equation}
S=P^-_{00}+R_\theta,\quad
K=\frac{1}{S}\begin{bmatrix}P^-_{00}\\P^-_{10}\end{bmatrix},\quad
\begin{bmatrix}\hat\theta^+\\\hat b^+\end{bmatrix}
=\begin{bmatrix}\hat\theta^-\\\hat b^-\end{bmatrix}
+\begin{bmatrix}K_\theta\\K_b\end{bmatrix}(\theta_a-\hat\theta^-).
\end{equation}
$K_\theta$ is dimensionless; $K_b$ has units s$^{-1}$ and is normally
negative with this sign convention. A positive angle innovation should
reduce an estimated positive gyro bias that has been over-subtracted.

The initialization is $P_0=\operatorname{diag}(1,0.01)$,
$\sigma_g^2=0.04$, $q_b=0.0001$ and $R_\theta=4$ in degree-based units.
At $\Delta=0.002$ s, the first prediction yields
\begin{align}
P^-_{00}&=1.000000200000267,\\
P^-_{01}=P^-_{10}&=-0.0000200002,\\
P^-_{11}&=0.0100002,\\
S&=5.000000200000267,\\
K_\theta&\simeq0.200000032,\qquad K_b\simeq-0.00000400004\ \mathrm{s}^{-1}.
\end{align}
The numerical first gain is not the eventual gain; plots and end-of-run
values are generated from the actual C trace.

\section{Accelerometer observations and their limitation}
For calibrated specific force $f_B$ in units of $g$, the small-acceleration
gravity observations are
\begin{equation}
\phi_a=\operatorname{atan2}(f_y,f_z),\qquad
\theta_a=\operatorname{atan2}(-f_x,\sqrt{f_y^2+f_z^2}).
\end{equation}
Convert these to degrees before feeding the C filter. The accelerometer
measures specific force, not gravity alone:
$f_B=R^{\mathsf T}(a_W+ge_3)$ in SI units. During sustained translation,
these angle observations are biased. Fixed $R_\theta=4$ deg$^2$ is a
provisional robustness choice, not the variance implied directly by
0.005 $g$ sensor noise. Their nonlinear Jacobian would propagate raw
measurement covariance; motion adds further model error. This explains
why a mathematically valid covariance recursion alone does not prove that
the reported confidence intervals are calibrated in real flight.

\section{Altitude filter: state, measurements and expanded products}
Use $x=[h,v_z]^{\mathsf T}$ and the exact $F,G$ derived earlier. Corrected
world-up acceleration is the control input. The C state prediction is
\begin{equation}
\hat h^-=\hat h+\Delta\hat v_z+\tfrac12\Delta^2 a_z,\qquad
\hat v_z^-=\hat v_z+\Delta a_z.
\end{equation}
The covariance equations from $FPF^{\mathsf T}+Q_{\rm held}$ are
\begin{align}
P^-_{00}&=P_{00}+\Delta(P_{01}+P_{10})+\Delta^2P_{11}
 +\tfrac14\sigma_a^2\Delta^4,\\
P^-_{01}&=P_{01}+\Delta P_{11}+\tfrac12\sigma_a^2\Delta^3,\\
P^-_{10}&=P_{10}+\Delta P_{11}+\tfrac12\sigma_a^2\Delta^3,\\
P^-_{11}&=P_{11}+\sigma_a^2\Delta^2.
\end{align}
Barometric height has $z_h=Hx+v_h$, $H=[1,0]$, $R_h=0.0025$ m$^2$ in
SITL. The gains are
\begin{equation}
K_h=\frac{P^-_{00}}{P^-_{00}+R_h},\qquad
K_v=\frac{P^-_{10}}{P^-_{00}+R_h}.
\end{equation}
Here $K_v$ is positive: a higher observed altitude tends to raise both
estimated height and upward speed. A barometer reading is corrected
\emph{once when fresh}, every ten control steps. Reusing the same reading as
if it were ten independent measurements would understate uncertainty.
The gains logged between observations are the most recent gains, not ten
new corrections.

SITL starts at $P_0=\operatorname{diag}(0.04,0.1)$, with $\sigma_a=1.5$.
The first 2 ms prediction gives
\begin{align}
P^-&=\begin{bmatrix}0.040000400009&0.000200009\\
0.000200009&0.100009\end{bmatrix},\\
S&=0.042500400009,\\
K_h&\simeq0.941177024,\qquad K_v\simeq0.00470605\ \mathrm{s}^{-1}.
\end{align}
The firmware scaffold's boot defaults are instead $P_0=\operatorname{diag}(1,1)$
and $R_h=2$ m$^2$. The host simulation deliberately supplies the documented
test configuration through its reset interface. Neither is a measured BMP280
noise model. They must not be confused when reproducing results.

\section{Joseph multiplication in full detail}
Let a symmetric prior be $P^-=\begin{bmatrix}a&b\\b&d\end{bmatrix}$,
measurement variance $r$, and gain $K=[k_0,k_1]^{\mathsf T}$.
For $H=[1,0]$,
\begin{equation}
M=I-KH=\begin{bmatrix}1-k_0&0\\-k_1&1\end{bmatrix}.
\end{equation}
Writing $z=1-k_0$ and multiplying row by column,
\begin{align}
MP^-&=\begin{bmatrix}za&zb\\b-k_1a&d-k_1b\end{bmatrix},\\
MP^-M^{\mathsf T}&=\begin{bmatrix}
z^2a&z(b-k_1a)\\z(b-k_1a)&d-2k_1b+k_1^2a
\end{bmatrix},\\
KrK^{\mathsf T}&=r\begin{bmatrix}k_0^2&k_0k_1\\k_0k_1&k_1^2\end{bmatrix}.
\end{align}
Adding gives exactly the four expressions implemented in
\texttt{covariance\_correct}. The routine averages the two off-diagonal
prior entries to remove small numerical asymmetry, and writes the same
posterior off-diagonal value twice. Joseph form is a sum of positive
semidefinite matrices in exact arithmetic. Finite-precision roundoff still
requires monitoring; a failed covariance validity check must not be disguised
by claiming the physical filter is healthy.

\section{Gravity compensation, lidar and observability}
The upward acceleration supplied to the filter is
\begin{equation}
a_z=-\sin\hat\theta\,f_x
 +\sin\hat\phi\cos\hat\theta\,f_y
 +\cos\hat\phi\cos\hat\theta\,f_z-g,
\end{equation}
with $f$ in m/s$^2$. At rest and level it is zero. In the CPU model this
uses the preceding C attitude estimate, not the exact plant attitude, so
attitude error affects height estimation. The resulting correlated errors
are not fully represented by this two-state covariance model.

For flat ground and a truly downward aligned rangefinder, approximate
$h=d\cos\phi\cos\theta$ after accounting for the sensor lever arm. Its
variance can be approximated by $R_h=J_d\Sigma J_d^{\mathsf T}$, with
\begin{equation}
J_d=\begin{bmatrix}\cos\phi\cos\theta&-d\sin\phi\cos\theta&-d\cos\phi\sin\theta\end{bmatrix}.
\end{equation}
Terrain slope, missing returns, tilt and range limits violate this simple
measurement model. A downward TOFSense-F2 Mini is therefore a useful future
height aid, but it is not fused in the current trace. A front sensor sees
only one ray; it cannot establish a full safe corridor by itself.

For the altitude pair, the two-observation observability matrix is
\begin{equation}
\mathcal O=\begin{bmatrix}H\\HF\end{bmatrix}
=\begin{bmatrix}1&0\\1&\Delta\end{bmatrix},\quad
\det\mathcal O=\Delta\ne0.
\end{equation}
Thus repeated independent height measurements make both height and speed
observable. The analogous angle-bias determinant is $-\Delta$. With all
height measurements missing, integration alone drifts and no longer bounds
height uncertainty. Loss of sensor aiding requires a defined, tested degraded
mode; generic zero motor output in flight is not a safe landing policy.

\chapter{Uncertainty, statistical checks and error accounting}
\section{What standard deviation does and does not mean}
For a scalar zero-mean Gaussian estimation error $e\sim\mathcal N(0,P)$,
$\sigma=\sqrt P$. The standardized variable $u=e/\sigma$ has unit normal
density $\exp(-u^2/2)/\sqrt{2\pi}$. Therefore
\begin{equation}
\Pr(|e|\leq c\sigma)=\int_{-c}^c\frac{e^{-u^2/2}}{\sqrt{2\pi}}\dd u,
\end{equation}
which is approximately 68.27\% for $c=1$ and 95.45\% for $c=2$.
A covariance is a model prediction, not a measured error bound. Gaussian
intervals can fail under bias, bad calibration, outliers, latency or incorrect
process-noise assumptions.

For $N$ Monte Carlo errors $e_i$ at a fixed time, compute
\begin{equation}
\bar e=\frac1N\sum_i e_i,\quad
s_e=\sqrt{\frac{\sum_i(e_i-\bar e)^2}{N-1}},\quad
\mathrm{RMSE}=\sqrt{\frac1N\sum_i e_i^2}.
\end{equation}
The $N-1$ denominator makes sample variance unbiased for independent samples.
RMSE includes bias; standard deviation describes spread about the sample mean.
For $N=500$ and true coverage 0.9545, the binomial standard error of an
observed two-sigma coverage fraction is
$\sqrt{0.9545(1-0.9545)/500}\simeq0.00932$ at one independent time point.
Successive time points within a trajectory are correlated, so pooling them
does not create that many independent trials.

\section{Innovation consistency and physical-parameter errors}
For a scalar observation, normalized innovation squared is
$\mathrm{NIS}=\nu^2/S$. Under the correctly specified linear Gaussian model,
its mean is one and its distribution is $\chi^2_1$. Persistent large values
suggest underestimated uncertainty, bias, bad timing or sensor failures;
very small values can indicate overestimated noise. This is a diagnostic,
not an automatic proof of sensor validity. The C trace records angle
innovation and its variance. Altitude innovation variance can be reconstructed
on fresh corrections from $S=R_h/(1-K_h)$, subject to roundoff and $K_h<1$.

For a physical parameter vector $\eta$ with covariance $\Sigma_\eta$, a
first-order output uncertainty is
\begin{equation}
\Sigma_f\simeq J_f\Sigma_\eta J_f^{\mathsf T},\qquad
(J_f)_{ij}=\frac{\partial f_i}{\partial\eta_j}.
\end{equation}
For example, $c_0=mg/(4T_{\max})$ has derivatives
\begin{equation}
\frac{\partial c_0}{\partial m}=\frac{g}{4T_{\max}},\qquad
\frac{\partial c_0}{\partial T_{\max}}=-\frac{mg}{4T_{\max}^2}.
\end{equation}
With independent
relative errors this gives
\begin{equation}
\left(\frac{\sigma_{c_0}}{c_0}\right)^2
\simeq\left(\frac{\sigma_m}{m}\right)^2
+\left(\frac{\sigma_T}{T_{\max}}\right)^2.
\end{equation}
For $b_x\propto aT_{\max}/J_x$, the logarithmic first-order perturbation is
$\delta b_x/b_x\simeq\delta a/a+\delta T/T-\delta J_x/J_x$.
Correlations must be included if, for example, moving a battery changes both
centre of mass and inertia. The current mass/thrust perturbations are scenario
tests, not measured probability distributions for the actual aircraft.

\section{Separate error sources}
The report distinguishes: floating-point and timestep error; noisy sensors;
calibration error; structural model error; uncertain mass/inertia/thrust;
latency and stale-data error; and perception/localization error. Running the
same approximate equations very precisely only reduces numerical error. It
does not eliminate structural or identification error. Six successful
scenarios provide evidence for those scenarios, not a probability of safe
flight in arbitrary environments.

\chapter{Companion navigation and perception}
\section{A bounded static-map baseline}
The common C++ planner discretizes $x\in[-2,10]$, $y\in[-4,4]$ into 0.2 m
cells. A circle obstacle of radius $r_i$ blocks centres within $r_i+0.65$ m.
The 0.65 m inflation includes the 0.34 m propeller envelope and additional
tracking/grid allowance. Eight-connected A* uses edge length $0.2$ or
$0.2\sqrt2$ and Euclidean heuristic. Diagonal motion through blocked corners
is rejected. Since Euclidean straight-line distance never exceeds path length,
the heuristic is admissible. It does not prove continuous-space safety under
arbitrary tracking errors; actual trajectory clearance is computed separately.

For waypoint error $e_{xy}$, set desired horizontal velocity
$v_c=\operatorname{clip}_{\|\cdot\|\leq0.8}(e_{xy})$ and acceleration
$a_{xy,c}=\operatorname{clip}_{\rm component,1.4}[1.8(v_c-v_{xy})]$.
With yaw $\psi$, transform world horizontal acceleration into body heading:
\begin{equation}
\begin{bmatrix}a_{xB}\\a_{yB}\end{bmatrix}
=\begin{bmatrix}\cos\psi&\sin\psi\\-\sin\psi&\cos\psi\end{bmatrix}
\begin{bmatrix}a_{xW}\\a_{yW}\end{bmatrix}.
\end{equation}
The approximate requests are $\theta_c=\arctan(a_{xB}/g)$ and
$\phi_c=-\arctan(a_{yB}/g)$, limited to 15 degrees. The waypoint advances
within 0.4 m. Goal acceptance requires horizontal error below 0.2 m,
vertical error below 0.2 m and horizontal speed below 0.15 m/s, followed by
one second dwell in the route test. Flights begin already airborne at 1.5 m;
these tests do not include takeoff, landing or recovery from a crash.

\section{ROS 2 data contract}
The ROS C++ source uses a custom timestamped AutonomySetpoint message with
roll angle, pitch angle, yaw rate, centered throttle, validity and goal status.
Angles are radians at the ROS boundary and degrees inside the C core. A
250 ms odometry freshness gate invalidates companion output. Local pose and
velocity must use one consistent frame; the wrapper requires a documented
adapter that transforms standard body-frame Odometry twist to world velocity.
Ordinary body twist must not be silently interpreted as world velocity.
No setpoint arms hardware. A future physical bridge must enforce pilot mode
authority, timestamp/sequence checks, watchdogs and a tested degraded mode.

The executed DDS test uses two Python adapter nodes with numeric test
envelopes, carrying sensor samples, commands, motor proxies and estimator
diagnostics. It exercises actual DDS serialization/transport and C execution,
but does not validate generated custom-message C++ bindings, physical serial
framing or a complete real-time companion system.

\section{Why YOLO is not a range sensor}
For a calibrated pinhole camera, optical coordinates satisfy
$u=f_xX/Z+c_x$ and $v=f_yY/Z+c_y$. If a registered depth image supplies $Z$,
\begin{equation}
X=(u-c_x)Z/f_x,\qquad Y=(v-c_y)Z/f_y.
\end{equation}
Without depth or another geometric constraint, multiplying $(X,Y,Z)$ by any
positive scale leaves $(u,v)$ unchanged. This proves the metric-scale
ambiguity of a single detection image. A bounding box cannot supply reliable
distance by itself. The tested geometry utility rejects invalid range values
and back-projects \emph{all} valid sampled depth pixels, including objects
outside YOLO's recognized classes. Bounding-box median depth is only a
semantic annotation.

Transform each point with calibrated extrinsics and synchronized pose:
$p_W=R_{WB}(R_{BC}p_C+t_{BC})+t_{WB}$. Future map insertion needs ground
filtering, voxelization, obstacle growth, clearing/decay, temporal consistency
and replanning. None of those steps can be replaced by the standalone bus-image
detection test. The same ROS C++ navigation core can run on Pi or Jetson;
perception acceleration and camera drivers remain platform-specific concerns.

\chapter{Reproduction and next physical measurements}
The evidence is dated 20 October 2026. Run the host C checks, build SITL
libraries, execute the route suite, execute the ROS DDS adapter test, and render
graphs/video using the scripts listed in the results chapter. All stochastic
seeds and assumed parameters are saved alongside the traces. The video plays
one simulated second per second; producing it offline is not proof of live
500 Hz execution. DDS wall-clock percentiles include Windows scheduling.

To make the model physically precise, measure complete takeoff mass and
centre of mass; identify motor/propeller thrust and current versus command at
several battery voltages; estimate actuator lag; measure inertia or assemble
component inertias; measure vibration, gyro Allan behaviour, accelerometer
and barometer noise after mounting; measure communication age/jitter; calibrate
camera intrinsics/extrinsics and rangefinder offsets. Until then all numeric
PID values in this report are simulation candidates. The LiPo's capacity and
C-rating alone do not establish usable energy, voltage sag, thrust reserve or
safe cell condition.

Primary implementation references for later integration are NVIDIA's
\href{https://docs.isaacsim.omniverse.nvidia.com/6.0.1/installation/requirements.html}{Isaac Sim requirements}
and Ultralytics' \href{https://docs.ultralytics.com/modes/predict/}{prediction}
and \href{https://docs.ultralytics.com/modes/export/}{model export} documentation.
The algebra and gain derivations above are developed explicitly for this project.
